% ============================================================
% 05_future_work.tex — Идна работа (македонски)
% ============================================================

\section{Дополнителни карактеристики на играчите}

\subsection{Карактеристики за замор и тренд на минути}

Тековниот модел не вклучува директни мерки за заморот на играчите или за трендовите на нивното одиграно време. Карактеристиките \texttt{minutes\_last3} и \texttt{minutes\_last5} (тркалачките просеци на одиграните минути во последните три и пет натпреварувачки недели, соодветно) би служеле како прокси за тековниот физички товар на играчот. Играч кој конзистентно одиграл по 90 и повеќе минути низ пет последователни натпреварувачки недели се соочува со значајно зголемен ризик од замор во следниот натпревар, особено кога тој натпревар се одвива во рамките на густ распоред од два или три натпревари неделно.

Имплементацијата на овие карактеристики би барала измени во pipeline-от за инженеринг на карактеристики (\texttt{feature\_engineering\_stage6.py}) за да се пресметаат тркалачките прозорци на минутите, последовано од целосно преобучување на сите четири модели по позиција. Прелиминарната корелациска анализа покажува позитивна врска меѓу \texttt{minutes\_last3} и \texttt{total\_points}, особено кај играчите со неправилни распределби на минутите (на пример, кандидатите за ротација кои повремено стартуваат). Клучната корист е раздвојувањето на играчите кои се во добра форма но уморни од оние кои се во добра форма и целосно фит.

Деривативната карактеристика \texttt{minutes\_trend} — означениот наклон на минутите низ последните три натпревари — би открила играчи чие учество е во пораст или во опаѓање. Бранител кој одиграл 45, 60 и 80 минути во три последователни натпревари веројатно се враќа во полна форма и заслужува позитивно прилагодување, додека спротивната траекторија би оправдала прилагодување надолу.

\subsection{Карактеристики за историја на повреди}

Историските податоци за повреди — тип на повреда, времетраење и фреквенција на повторување — претставуваат потенцијално вреден сигнал кој моментално го нема во моделот. Играчите со документирана историја на повреди на хамстрингот или мускулни повреди носат зголемен ризик од повторување во периодите со високи физички барања. Иако оваа информација е достапна од трети страни како Transfermarkt и PhysioRoom, нејзината интеграција бара внимателно моделирање за да се избегне пристрасност при селекција (играчите со тешки истории на повреди едноставно може да ја напуштат лигата) и неконзистентност во етикетирањето меѓу изворите.

Практичната имплементација би одржувала регистар на повреди по играч, ажуриран од податоците на FPL API во intel\_01, при што секој настан на повреда би се следел според почетната натпреварувачка недела и времетраењето. Овој регистар би генерирал карактеристики како \texttt{days\_since\_last\_injury}, \texttt{injury\_count\_season} и \texttt{injury\_type\_severity}, кои би го дополниле моделот за ризик од ротација (intel\_04) со сигнал за повреди со подолга меморија.

\section{Активна имплементација во ГН30+}

\subsection{Поврзување со вистинскиот FPL API}

Системот, како што е евалуиран, работи во режим на симулација, повторно репродуцирајќи историски податоци по натпреварувачки недели. Следниот природен развоен чекор е активна имплементација од ГН30 натаму, со директно поврзување со FPL API за извршување на вистински промени во тимот. Тоа би барало:

\begin{enumerate}
  \item Автентикација со FPL API преку акредитиви на вистинска менаџерска сметка.
  \item Имплементација на крајната точка за поднесување трансфери, со испраќање на препорачаните трансфери на системот до серверите на FPL пред секој рок.
  \item Оперативно следење за потврда дека поднесените трансфери се прифатени и правилно евидентирани.
  \item Автоматизирано собирање на податоци по натпреварувачката недела: преземање на вистинските поени, ажурирање на тренинг множеството и активирање на целосниот pipeline за преобучување на моделите.
\end{enumerate}

Активната имплементација би обезбедила независна надворешна валидација на перформансите на системот: резултатите би биле директно споредливи со глобалните FPL рангирања (просек од приближно 50 поени по натпреварувачка недела во 2024--25) и со перформансите на топ 100K менаџерите, кои служат како практичен бенчмарк за врвно FPL менаџирање.

\subsection{Архитектура на деплојментот}

За продукциски деплојмент, pipeline-от би бил рефакторизиран во микросервисна архитектура распоредена преку cron задачи. Intel pipeline-от (модули 01--05) би се извршувал приближно три часа пред секој рок, ILP оптимизаторот би се извршувал триесет минути пред рокот со финално прилагодените предикции, а резултатите би се прикажувале преку постоечкиот Flask UI. Ракувањето со грешки би вклучувало враќање на составот од минатата недела во случај ILP да не успее да произведе валидно решение во рамките на временскиот прозорец.

\section{Пристап со учење со засилување}

\subsection{Формулација на MDP}

Управувањето со FPL низ цела сезона природно може да се претстави како Марков Одлучен Процес (MDP), што го прави кандидат за пристапи со учење со засилување:

\begin{itemize}
  \item \textbf{Состојба} $s_t$: тековниот состав на тимот, преостанатиот буџет, банкираните слободни трансфери, преостанатите чипови, профилите на формата на играчите и распоредот на FDR за следните $N$ натпреварувачки недели.
  \item \textbf{Акција} $a_t$: множество одлуки за трансфери (купување/продавање поединечни играчи), избор на капитен и опционо активирање чип.
  \item \textbf{Награда} $r_t$: вистинските поени освоени во натпреварувачката недела $t$, минус казнените поени од трансфери.
  \item \textbf{Транзиции}: стохастички, зависно од непредвидливите реални перформанси на играчите.
\end{itemize}

Алгоритми како Deep Q-Networks (DQN) или Proximal Policy Optimization (PPO) би можеле директно да го оптимизираат долгорочниот кумулативен принос, наместо предвидените поени за една натпреварувачка недела што ги максимизира тековниот ILP. Тоа е особено корисно за одлуките за чипови, кои носат импликации за повеќе натпреварувачки недели: чипот Bench Boost е најкорисен кога се применува во двојна натпреварувачка недела во која играчите од клупата имаат добри натпревари, а тоа може да е оддалечено неколку недели од тековната точка на одлучување.

\subsection{Клучни предизвици}

Главниот предизвик за пристапот со учење со засилување е огромниот простор на состојба и акција: над 800 достапни играчи, 15 позиции во составот, континуирана динамика на цените и сезона со ограничен хоризонт од 38 натпреварувачки недели. Ефикасното тренирање бара симулирани FPL средини изградени од историските сезонски податоци, во кои стохастичките исходи на перформансите можат да се земаат како примероци од емпириската распределба на поените на играчите условена од нивната форма и натпреварот. Постоечкиот сезонски симулатор е природна почетна точка за конструирање на таква средина.

Пристапите за учење со засилување базирани на модел, кои најпрво учат модел на свет (предвидувајќи распределби на поени) и потоа планираат во него, можат да бидат особено добро прилагодени, бидејќи моделите за ML предикција веќе ја обезбедуваат основната компонента на моделот на свет.

\section{Подобрување на покриеноста на прес-конференциите}

\subsection{Поправање на ограничувањето за Newcastle}

Документираното ограничување на intel\_02 — пропуштањето клубови кои не се спомнати во насловите на написите — може да се адресира преку неколку комплементарни пристапи:

\begin{itemize}
  \item \textbf{Совпаѓање на имиња меѓу клубови}: Пребарување на сите познати имиња на играчи низ целиот текст на написите, без оглед на тоа под чија клупска секција паѓа текстот. Тоа беше тестирано но произведе нето намалување на резултатот од приближно 113 поени поради каскадни промени во составот. Повнимателно калибриран праг на сигурност за совпаѓањата меѓу клубови може да го избегне ваквото назадување.
  \item \textbf{Стримови на официјалните клупски прес-конференции}: Скрапирање на официјалните клупски веб-сајтови или YouTube канали за транскрипти од прес-конференциите по тренинзите, кои обично се објавуваат 48 часа пред секој натпревар.
  \item \textbf{Интеграција со Twitter/X}: Следење во реално време на сметките на новинарите кои покриваат специфични клубови ги дава најнавремените вести за повреди, често објавени само неколку минути по тренингот.
\end{itemize}

\subsection{Спојување на повеќе извори на intel}

Иден модул intel\_08 би можел да имплементира принципиелна рамка за спојување на информации од повеќе извори. На секој извор (FPL API, Fantasy Football Scout, Premier Injuries, клупски веб-сајтови, социјални мрежи) би му се доделил пондер на доверливост врз основа на неговата историска точност во предвидувањето на вистинските исходи на статусот на играчот. Бајесово правило за ажурирање би ги комбинирало сигналите од повеќе извори во апостериорна проценка на достапноста:

\begin{equation}
P(\text{достапен} \mid \text{сигнали}) \propto P(\text{достапен}) \prod_k P(\text{сигнал}_k \mid \text{достапен})^{w_k}
\end{equation}

каде што $w_k$ е пондерот на доверливост за изворот $k$, калибриран од историските податоци.

\section{Спречување на продај-купи назад}

\subsection{Проблемот}

Тековниот ILP нема меморија за одлуките за трансфери од минатата недела, што дозволува сценарија во кои истиот играч е продаден во една натпреварувачка недела и повторно купен во следната. Иако тоа е технички дозволено според правилата на FPL (по цена од четири поени по трансфер), таквата шема укажува или на нестабилен модел за предикција или на субоптимална трансферна стратегија која би можела да се исправи преку механизам со меморија.

\subsection{Можни решенија}

Член за казна при повторно купување, додаден на целната функција на ILP, го обесхрабрува повторното купување на неодамна продаден играч. Оптималната големина на казната зависи од очекуваната корист од повторното купување: ако предвидениот резултат на играчот во неделата на повторното купување значително ја надминува казната, повторното купување е рационално. Динамичка казна која експоненцијално опаѓа со бројот на натпреварувачки недели од продажбата — висока во непосредно следната натпреварувачка недела, приближувајќи се кон нула по пет или повеќе натпреварувачки недели — би постигнала подобар баланс од фиксна казна, истовремено избегнувајќи го назадувањето на резултатот забележано кога беше тестирана рамна казна. Калибрирањето на оваа стапка на опаѓање е и самото хиперпараметар кој би можел да се вклучи во идно Optuna пребарување.
